\section{AHSME 1978}
        \begin{enumerate}
        	\item (\textit{AHSME 1978}) If $1-\frac{4}{x}+\frac{4}{x^2}=0$, then $\frac{2}{x}$ equals


 $\textbf{(A) }-1\qquad \textbf{(B) }1\qquad \textbf{(C) }2\qquad \textbf{(D) }-1\text{ or }2\qquad  \textbf{(E) }-1\text{ or }-2$


 

	\item (\textit{AHSME 1978}) If four times the reciprocal of the circumference of a circle equals the diameter of the circle, then the area of the circle is 


 $\textbf{(A) }\frac{1}{\pi^2}\qquad \textbf{(B) }\frac{1}{\pi}\qquad \textbf{(C) }1\qquad \textbf{(D) }\pi\qquad \textbf{(E) }\pi^2$


 

	\item (\textit{AHSME 1978}) For all non-zero numbers $x$ and $y$ such that $x = 1/y$, $\left(x-\frac{1}{x}\right)\left(y+\frac{1}{y}\right)$ equals


 $\textbf{(A) }2x^2\qquad \textbf{(B) }2y^2\qquad \textbf{(C) }x^2+y^2\qquad \textbf{(D) }x^2-y^2\qquad  \textbf{(E) }y^2-x^2$


 

	\item (\textit{AHSME 1978}) If $a = 1,~ b = 10, ~c = 100$, and $d = 1000$, then $(a+ b+ c-d) + (a + b- c+ d) +(a-b+ c+d)+ (-a+ b+c+d)$ is equal to


 $\textbf{(A) }1111\qquad \textbf{(B) }2222\qquad \textbf{(C) }3333\qquad \textbf{(D) }1212\qquad  \textbf{(E) }4242$


 

	\item (\textit{AHSME 1978}) Four boys bought a boat for $\textdollar 60$. The first boy paid one half of the sum of the amounts paid by the other boys; 
the second boy paid one third of the sum of the amounts paid by the other boys; 
and the third boy paid one fourth of the sum of the amounts paid by the other boys. How much did the fourth boy pay? 


 $\textbf{(A) }\textdollar 10\qquad \textbf{(B) }\textdollar 12\qquad \textbf{(C) }\textdollar 13\qquad \textbf{(D) }\textdollar 14\qquad \textbf{(E) }\textdollar 15$


 

	\item (\textit{AHSME 1978}) The number of distinct pairs $(x,y)$ of real numbers satisfying both of the following equations: 


 
 \[x=x^2+y^2 \    \ y=2xy\] 
is


 $\textbf{(A) }0\qquad \textbf{(B) }1\qquad \textbf{(C) }2\qquad \textbf{(D) }3\qquad  \textbf{(E) }4$


 

	\item (\textit{AHSME 1978}) Opposite sides of a regular hexagon are $12$ inches apart. The length of each side, in inches, is


 $\textbf{(A) }7.5\qquad \textbf{(B) }6\sqrt{2}\qquad \textbf{(C) }5\sqrt{2}\qquad \textbf{(D) }\frac{9}{2}\sqrt{3}\qquad \textbf{(E) }4\sqrt{3}$


 

	\item (\textit{AHSME 1978}) If $x\neq y$ and the sequences $x,a_1,a_2,y$ and $x,b_1,b_2,b_3,y$ each are in arithmetic progression, then $(a_2-a_1)/(b_2-b_1)$ equals


 $\textbf{(A) }\frac{2}{3}\qquad \textbf{(B) }\frac{3}{4}\qquad \textbf{(C) }1\qquad \textbf{(D) }\frac{4}{3}\qquad \textbf{(E) }\frac{3}{2}$


 

	\item (\textit{AHSME 1978}) If $x<0$, then $\left|x-\sqrt{(x-1)^2}\right|$ equals


 $\textbf{(A) }1\qquad \textbf{(B) }1-2x\qquad \textbf{(C) }-2x-1\qquad \textbf{(D) }1+2x\qquad  \textbf{(E) }2x-1$


 

	\item (\textit{AHSME 1978}) If $\mathit{B}$ is a point on circle $\mathit{C}$ with center $\mathit{P}$, then the set of all points $\mathit{A}$ in the plane of circle 
$\mathit{C}$ such that the distance between $\mathit{A}$ and $\mathit{B}$ is less than or equal to the distance between $\mathit{A}$ 
and any other point on circle $\mathit{C}$ is


 $\textbf{(A) }\text{the line segment from }P \text{ to }B\qquad\\ \textbf{(B) }\text{the ray beginning at }P \text{ and passing through }B\qquad\\ \textbf{(C) }\text{a ray beginning at }B\qquad\\ \textbf{(D) }\text{a circle whose center is }P\qquad\\ \textbf{(E) }\text{a circle whose center is }B$


 

	\item (\textit{AHSME 1978}) If $r$ is positive and the line whose equation is $x + y = r$ is tangent to the circle whose equation is $x^2 + y ^2 = r$, then $r$ equals


 $\textbf{(A) }\frac{1}{2}\qquad \textbf{(B) }1\qquad \textbf{(C) }2\qquad \textbf{(D) }\sqrt{2}\qquad  \textbf{(E) }2\sqrt{2}$


 

	\item (\textit{AHSME 1978}) In $\triangle ADE$, $\measuredangle ADE=140^\circ$, points $B$ and $C$ lie on sides $AD$ and $AE$, 
respectively, and points $A,~B,~C,~D,~E$ are distinct.* If lengths $AB,~BC,~CD$, and $DE$ are all equal, 
then the measure of $\measuredangle EAD$ is


 
                \begin{itemize}
                    \item The specification that points $A,B,C,D,E$ be distinct was not included in the original statement of the problem.
                \end{itemize}


                If $B=D$, then $C=E$ and $\measuredangle EAD=20^\circ$.     


 $\textbf{(A) }5^\circ\qquad \textbf{(B) }6^\circ\qquad \textbf{(C) }7.5^\circ\qquad \textbf{(D) }8^\circ\qquad \textbf{(E) }10^\circ$


 

	\item (\textit{AHSME 1978}) If $a,b,c$, and $d$ are non-zero numbers such that $c$ and $d$ are the solutions of $x^2+ax+b=0$ and $a$ and $b$ are 
the solutions of $x^2+cx+d=0$, then $a+b+c+d$ equals


 $\textbf{(A) }0\qquad \textbf{(B) }-2\qquad \textbf{(C) }2\qquad \textbf{(D) }4\qquad  \textbf{(E) }(-1+\sqrt{5})/2$


 

	\item (\textit{AHSME 1978}) If an integer $n > 8$ is a solution of the equation $x^2 - ax+b=0$ and the representation of $a$ in the base-$n$ number system is $18$, 
then the base-n representation of $b$ is


 $\textbf{(A)}\ 18 \qquad \textbf{(B)}\ 20 \qquad \textbf{(C)}\ 80 \qquad \textbf{(D)}\ 81 \qquad \textbf{(E)}\ 280$


 

	\item (\textit{AHSME 1978}) If $\sin x+\cos x=1/5$ and $0\le x<\pi$, then $\tan x$ is


 $\textbf{(A) }-\frac{4}{3}\qquad \textbf{(B) }-\frac{3}{4}\qquad \textbf{(C) }\frac{3}{4}\qquad \textbf{(D) }\frac{4}{3}\qquad\\ \textbf{(E) }\text{not completely determined by the given information}$


 

	\item (\textit{AHSME 1978}) In a room containing $N$ people, $N > 3$, at least one person has not shaken hands with everyone else in the room. 
What is the maximum number of people in the room that could have shaken hands with everyone else?


 $\textbf{(A) }0\qquad \textbf{(B) }1\qquad \textbf{(C) }N-1\qquad \textbf{(D) }N\qquad  \textbf{(E) }\text{none of these}$


 

	\item (\textit{AHSME 1978}) If $k$ is a positive number and $f$ is a function such that, for every positive number $x$, $\left[f(x^2+1)\right]^{\sqrt{x}}=k$; 
then, for every positive number $y$, $\left[f(\frac{9+y^2}{y^2})\right]^{\sqrt{\frac{12}{y}}}$ is equal to


 $\textbf{(A) }\sqrt{k}\qquad \textbf{(B) }2k\qquad \textbf{(C) }k\sqrt{k}\qquad \textbf{(D) }k^2\qquad  \textbf{(E) }y\sqrt{k}$


 

	\item (\textit{AHSME 1978}) What is the smallest positive integer $n$ such that $\sqrt{n}-\sqrt{n-1}<.01$?


 $\textbf{(A) }2499\qquad \textbf{(B) }2500\qquad \textbf{(C) }2501\qquad \textbf{(D) }10,000\qquad  \textbf{(E) }\text{There is no such integer}$


 

	\item (\textit{AHSME 1978}) A positive integer $n$ not exceeding $100$ is chosen in such a way that if $n\le 50$, then the probability of choosing 
$n$ is $p$, and if $n > 50$, then the probability of choosing $n$ is $3p$. The probability that a perfect square is chosen is


 $\textbf{(A) }.05\qquad \textbf{(B) }.065\qquad \textbf{(C) }.08\qquad \textbf{(D) }.09\qquad  \textbf{(E) }.1$


 

	\item (\textit{AHSME 1978}) If $a,b,c$ are non-zero real numbers such that 
 \[\frac{a+b-c}{c}=\frac{a-b+c}{b}=\frac{-a+b+c}{a},\] and 
 \[x=\frac{(a+b)(b+c)(c+a)}{abc},\] and $x<0,$ then $x$ equals


 $\textbf{(A) }{-}1\qquad \textbf{(B) }{-}2\qquad \textbf{(C) }{-}4\qquad \textbf{(D) }{-}6\qquad  \textbf{(E) }{-}8$


 

	\item (\textit{AHSME 1978}) For all positive numbers $x$ distinct from $1$,


 
 \[\frac{1}{\log_3(x)}+\frac{1}{\log_4(x)}+\frac{1}{\log_5(x)}\]


 equals


 $\text{(A) }\frac{1}{\log_{60}(x)}\qquad\\ \text{(B) }\frac{1}{\log_{x}(60)}\qquad\\ \text{(C) }\frac{1}{(\log_{3}(x))(\log_{4}(x))(\log_{5}(x))}\qquad\\ \text{(D) }\frac{12}{\log_{3}(x)+\log_{4}(x)+\log_{5}(x)}\qquad\\ \text{(E) }\frac{\log_{2}(x)}{\log_{3}(x)\log_{5}(x)}+\frac{\log_{3}(x)}{\log_{2}(x)\log_{5}(x)}+\frac{\log_{5}(x)}{\log_{2}(x)\log_{3}(x)}$


 

	\item (\textit{AHSME 1978}) The following four statements, and only these are found on a card:

                    \begin{figure}[H]
                        \centering
                        \includegraphics[width=0.5\textwidth]{Problem/AHSME_I_1978_22_Problem_Figure_1.png}
                    \end{figure}
                    


 (Assume each statement is either true or false.) Among them the number of false statements is exactly


 $\textbf{(A)}\ 0 \qquad \textbf{(B)}\ 1 \qquad \textbf{(C)}\ 2 \qquad \textbf{(D)}\ 3 \qquad \textbf{(E)}\ 4$


 

	\item (\textit{AHSME 1978}) 
                    \begin{figure}[H]
                        \centering
                        \includegraphics[width=0.5\textwidth]{Problem/AHSME_I_1978_23_Problem_Figure_1.png}
                    \end{figure}
                    


 Vertex $E$ of equilateral $\triangle ABE$ is in the interior of square $ABCD$, and $F$ is the point of intersection 
of diagonal $BD$ and line segment $AE$. If length $AB$ is $\sqrt{1+\sqrt{3}}$ then the area of $\triangle ABF$ is


 $\textbf{(A) }1\qquad \textbf{(B) }\frac{\sqrt{2}}{2}\qquad \textbf{(C) }\frac{\sqrt{3}}{2}\qquad \textbf{(D) }4-2\sqrt{3}\qquad  \textbf{(E) }\frac{1}{2}+\frac{\sqrt{3}}{4}$


 

	\item (\textit{AHSME 1978}) If the distinct non-zero numbers $x ( y - z),~ y(z - x),~ z(x - y )$ form a geometric progression with common ratio $r$, 
then $r$ satisfies the equation


 $\textbf{(A) }r^2+r+1=0\qquad \textbf{(B) }r^2-r+1=0\qquad \textbf{(C) }r^4+r^2-1=0\qquad\\ \textbf{(D) }(r+1)^4+r=0\qquad  \textbf{(E) }(r-1)^4+r=0$


 

	\item (\textit{AHSME 1978}) Let $a$ be a positive number. Consider the set $S$ of all points whose rectangular coordinates $(x, y )$ satisfy all of the following conditions:


 $\text{(i) }\frac{a}{2}\le x\le 2a\qquad \text{(ii) }\frac{a}{2}\le y\le 2a\qquad \text{(iii) }x+y\ge a\\ \\ \qquad \text{(iv) }x+a\ge y\qquad \text{(v) }y+a\ge x$


 The boundary of set S is a polygon with


 $\textbf{(A) }3\text{ sides}\qquad \textbf{(B) }4\text{ sides}\qquad \textbf{(C) }5\text{ sides}\qquad \textbf{(D) }6\text{ sides}\qquad \textbf{(E) }7\text{ sides}$


 

	\item (\textit{AHSME 1978}) 
                    \begin{figure}[H]
                        \centering
                        \includegraphics[width=0.5\textwidth]{Problem/AHSME_I_1978_26_Problem_Figure_1.png}
                    \end{figure}
                    


 

In $\triangle ABC, AB = 10~ AC = 8$ and $BC = 6$. Circle $P$ is the circle with smallest radius which passes through $C$ 
and is tangent to $AB$. Let $Q$ and $R$ be the points of intersection, distinct from $C$ , of circle $P$ with sides $AC$ 
and $BC$, respectively. The length of segment $QR$ is


 $\textbf{(A) }4.75\qquad \textbf{(B) }4.8\qquad \textbf{(C) }5\qquad \textbf{(D) }4\sqrt{2}\qquad  \textbf{(E) }3\sqrt{3}$


 

	\item (\textit{AHSME 1978}) There is more than one integer greater than $1$ which, when divided by any integer $k$ such that $2 \le k \le 11$, has a remainder of $1$. 
What is the difference between the two smallest such integers?


 $\textbf{(A) }2310\qquad \textbf{(B) }2311\qquad \textbf{(C) }27,720\qquad \textbf{(D) }27,721\qquad  \textbf{(E) }\text{none of these}$


 

	\item (\textit{AHSME 1978}) 
                    \begin{figure}[H]
                        \centering
                        \includegraphics[width=0.5\textwidth]{Problem/AHSME_I_1978_28_Problem_Figure_1.png}
                    \end{figure}
                    


 

If $\triangle A_1A_2A_3$ is equilateral and $A_{n+3}$ is the midpoint of line segment $A_nA_{n+1}$ for all positive integers $n$, 
then the measure of $\measuredangle A_{44}A_{45}A_{43}$ equals


 $\textbf{(A) }30^\circ\qquad \textbf{(B) }45^\circ\qquad \textbf{(C) }60^\circ\qquad \textbf{(D) }90^\circ\qquad  \textbf{(E) }120^\circ$


 

	\item (\textit{AHSME 1978}) Sides $AB,~ BC, ~CD$ and $DA$, respectively, of convex quadrilateral $ABCD$ are extended past $B,~ C ,~ D$ and $A$ to points $B',~C',~ D'$ and $A'$. 
Also, $AB = BB' = 6,~ BC = CC' = 7, ~CD = DD' = 8$ and $DA = AA' = 9$; and the area of $ABCD$ is $10$. The area of $A 'B 'C'D'$ is


 $\textbf{(A) }20\qquad \textbf{(B) }40\qquad \textbf{(C) }45\qquad \textbf{(D) }50\qquad  \textbf{(E) }60$


 

	\item (\textit{AHSME 1978}) In a tennis tournament, $n$ women and $2n$ men play, and each player plays exactly one match with every other player. 
If there are no ties and the ratio of the number of matches won by women to the number of matches won by men is $7/5$, then $n$ equals


 $\textbf{(A) }2\qquad \textbf{(B) }4\qquad \textbf{(C) }6\qquad \textbf{(D) }7\qquad  \textbf{(E) }\text{none of these}$


 

\end{enumerate}